\providecommand{\kBKM}{\kappa_{\mathrm{BKM}}}
\providecommand{\kcat}{\kappa_{\mathrm{cat}}}
\providecommand{\kch}{\kappa_{\mathrm{ch}}}
\providecommand{\kfib}{\kappa_{\mathrm{fiber}}}
\providecommand{\Mp}{\mathrm{Mp}}
\providecommand{\Sp}{\mathrm{Sp}}

\section*{Half-integral diagonal-divisor rows and super-VOA tests}

At the \(K3\times E\) boundary three automorphic parametrisations meet:
the compact CHL ladder, the Nikulin symplectic-order extension, and the
Gritsenko--Cl\'ery diagonal-divisor catalogue.

\[
\begin{array}{c|c|c}
\text{CHL order }N & c_N(0) & \kBKM(\Phi_N)=c_N(0)/2 \\
\hline
1&10&5\\
2&8&4\\
3&6&3\\
4&4&2\\
6&2&1
\end{array}
\]
These five rows are the compact \(K3\times E\) CHL CY\(_3\) ladder.
Since the source has dimension \(3\), the specialised native output of
\(\Phi_3\) is \(E_1\)-chiral.

\[
\begin{array}{c|c|c}
\text{Nikulin order }N & c_N(0) & \kBKM(\Phi_N)=c_N(0)/2 \\
\hline
5&4&2\\
7&2&1\\
8&2&1
\end{array}
\]
These three entries extend the same Borcherds-weight computation to
Nikulin symplectic orders beyond the compact CHL ladder.  Their weights
are integral: \(2,1,1\).  The elliptic-curve automorphism constraint
selects exactly \(N\in\{1,2,3,4,6\}\) for smooth \(K3\times E\) CHL
CY\(_3\) quotients; the rows \(N=5,7,8\) form the symplectic-order
extension outside the compact CHL host geometry.

The Gritsenko--Cl\'ery diagonal-divisor catalogue is indexed by triples
\((t,N;k)\):
\[
\begin{array}{c|c|c|c}
(t,N;k) & F_{t,N} & 2k & \text{multiplier}\\
\hline
(1,1;5)&\Delta_5&10&\chi_2\\
(2,1;2)&\Delta_2&4&\chi_4\\
(1,2;3)&\nabla_3&6&\chi_2\\
(3,1;1)&\Delta_1&2&\chi_6\\
(1,3;2)&\nabla_2&4&\chi_2\\
(4,1;\tfrac12)&\Delta_{1/2}&1&\chi_8\\
(1,4;\tfrac32)&\nabla_{3/2}&3&\chi_4\\
(2,2;1)&Q_1&2&\chi_4
\end{array}
\]
The half-integral entries are exactly
\[
  (4,1;\tfrac12),\qquad (1,4;\tfrac32).
\]
The twice-weight tuple of the catalogue is
\[
  (10,4,6,2,4,1,3,2).
\]
Half-integrality is a multiplier-system property of the
\((4,1;\tfrac12)\) and \((1,4;\tfrac32)\) diagonal-divisor rows.  The
Nikulin orders \(5,7,8\) above carry the integer weights \(2,1,1\).

\paragraph{Borcherds weight.}
For a Borcherds product with Jacobi input
\[
  \psi(\tau,z)=\sum c(n,\ell)q^n\zeta^\ell ,
\]
the scalar automorphy weight is read from the constant term:
\[
  \mathrm{wt}(\operatorname{Borch}(\psi))=\frac{c(0,0)}{2}.
\]
For the multi-cusp Gritsenko--Cl\'ery products the same statement is the
cusp-weighted zero-coefficient sum in their product theorem.  This
automorphic denominator weight is the invariant denoted \(\kBKM\).
Vertex-algebra central charge, Hodge supertrace, and product Euler
characteristic are separate invariants.

\paragraph{Denominator algebra versus super-VOA.}
A Borcherds product becomes a denominator function after the root datum
has been fixed: an integral Lorentzian lattice, a positive cone, a Weyl
vector, real simple roots, imaginary simple-root multiplicities, and a
parity rule.  A super-VOA requires further structure: a state-field
construction, a conformal vector, a local cocycle, an invariant form,
and a compatible odd sector.  A Friedan--Martinec--Shenker
\(\beta\gamma\)-\(bc\) system supplies such an odd sector with specified
fields and conformal weights.  Half-integral automorphy weight
supplies the multiplier system.  Spin modules, Neveu--Schwarz actions,
and vacuum super-VOAs require their own constructions.

\paragraph{Free-field character test.}
Consider the candidate rank-one FMS oscillator product
\[
\chi(q)=
q^{-c/24}
\prod_{n\geq 1}(1-q^{n-1/2})^{-1}(1-q^n)^{-2}.
\]
Writing \(t=q^{1/2}\), the truncation through \(t^4\) is finite:
\[
  \prod_{n\geq 1}(1-t^{2n-1})^{-1}(1-t^{2n})^{-2}
  =
  (1-t)^{-1}(1-t^2)^{-2}(1-t^3)^{-1}(1-t^4)^{-2}
  +O(t^5)
\]
and hence
\[
  \prod_{n\geq 1}(1-t^{2n-1})^{-1}(1-t^{2n})^{-2}
  =
  1+t+3t^2+4t^3+9t^4+\cdots .
\]
The spin-lattice theta character with initial terms
\[
  q^{-1/16}(1+q^{1/2}+q+2q^{3/2}+\cdots)
  =
  q^{-1/16}(1+t+t^2+2t^3+\cdots)
\]
fails this necessary coefficient test at \(t^2=q\), where the oscillator
product has coefficient \(3\) and the theta character has coefficient
\(1\).  The leading exponent \(q^{-c/24}=q^{-1/16}\) gives \(c=3/2\)
once a vacuum-energy normalisation is fixed; coefficient equality
is still required for a spin-lattice character.  The free-field product
is a test datum; construction of the required super-VOA requires the
additional state-field data above.

\paragraph{\(K3\times E\) boundary.}
The compact product \(Y=K3\times E\) carries five distinct readings:
\[
  \kcat(Y)=0,\qquad
  \kch(Y)=0,\qquad
  \kch^{\mathrm{Heis}}(Y)=3,\qquad
  \kBKM(\mathfrak g_{\Delta_5})=5,\qquad
  \kfib(Y)=24.
\]
Here \(\kcat(K3\times E)=\chi(\mathcal O_{K3})\chi(\mathcal O_E)=0\);
the compact Hodge supertrace \(\kch(K3\times E)=0\) by K\"unneth and
odd-dimensional Serre cancellation; \(\kch^{\mathrm{Heis}}(K3\times
E)=3\) is the Heisenberg--Mukai specialisation; \(\kBKM(\mathfrak
g_{\Delta_5})=5\) is the Borcherds weight of \(\Delta_5\); and
\(\kfib=24\) is the K3 Mukai-lattice rank.  The value
\(\chi(\mathcal O_{K3})=2\) belongs to the K3 fibre.  At \(d\geq 3\)
the native \(\Phi_d\)-output is \(E_1\)-chiral; any \(E_2\) structure
belongs to a derived centre, Drinfeld double, or representation category
after its own hypotheses are imposed.

The primitive K3 denominator normalisation is \(\Delta_5\), with
\(\kBKM(\Delta_5)=5\).  The monic product used in the reduced-DT scalar
normalisation is
\[
  D_5 = 64^{-1}\Delta_5,\qquad
  \Phi_{10}^{\mathrm{OP}} = D_5^2 = 4096^{-1}\Delta_5^2 .
\]
The unnormalised Igusa square \(\Phi_{10}^{\mathrm{un}}=\Delta_5^2\)
has weight \(10\); the primitive BKM denominator reading remains
\(\kBKM(\Delta_5)=5\).

\paragraph{References.}
Borcherds 1998, \emph{The singular theta correspondence and
automorphic forms}, Theorem~13.3; Gritsenko--Nikulin 1995 Part~II,
Theorem~2.1; Gritsenko 1999, \emph{Arithmetical lifting and its
applications}, Theorem~3.4; Gritsenko--Cl\'ery 2008,
\emph{The Siegel modular forms of genus \(2\) with the simplest
divisor}, Theorems~1.4 and~3.1; Nikulin 1980 on symplectic
automorphisms of K3 surfaces; Kac 1998, Chapters~4--5;
Frenkel--Lepowsky--Meurman 1988, Chapters~5--8;
Friedan--Martinec--Shenker 1986, first-order systems.
